\section{Database Setup and Loading}
\label{sec:dbsetup}

\subsection{Stack}
We use PostgreSQL 16 \cite{postgresql} running inside a single Docker
container defined in \texttt{docker-compose.yml} \cite{docker}. Schema
migrations are managed by Alembic \cite{alembic} so that any team
member can rebuild the database from a clean state with a single
command. Application code is Python 3.12 with SQLAlchemy ORM models
for parsing and loading.

\subsection{Repository Layout}
\begin{lstlisting}[basicstyle=\ttfamily\scriptsize, frame=none, language={}]
data-201-security-log-analysis/
|-- alembic/              schema migrations
|-- docker-compose.yml    Postgres 16
|-- src/
|   |-- models/           SQLAlchemy models
|   |-- parsers/          raw -> dict parsers
|   |-- loaders/          dict -> row loaders
|   `-- dashboard/        Streamlit app + db.py
|-- sql/
|   |-- 3nf/              per-table CREATE TABLE
|   |-- 3nf/_indexes.sql  canonical index DDL
|   |-- 3nf/_views.sql    canonical view DDL
|   |-- queries/basic/    6 basic .sql files
|   `-- queries/advanced/ 13 advanced .sql files
|-- notebooks/            10 exploration nbs
|-- docs/                 schema, ERD, findings
`-- report/               this report
\end{lstlisting}

\subsection{Reproducible Load}
The full load runs in five commands from a clean clone of the
repository:

\begin{lstlisting}[language=bash, basicstyle=\ttfamily\small, frame=single]
docker compose up -d                  # Postgres on :5432
python -m venv venv && source venv/bin/activate
pip install -r requirements.txt
alembic -c alembic/alembic.ini upgrade head
python -m src.loaders.run_all
\end{lstlisting}

The migration chain creates 22 tables across staging and 3NF schemas;
the loader script reads the raw \texttt{russellmitchell} files, runs
the parsers in \texttt{src/parsers/}, and inserts rows in dependency
order (host inventory first, then audit events, then attack labels).
The loader is idempotent: re-running against an already-loaded database
truncates and re-inserts each table inside a transaction, so the
output is reproducible regardless of starting state.

\subsection{Migration Head and Schema Objects}
After \texttt{alembic upgrade head} the database carries revision
\texttt{742e860d116f}, which adds two analytical objects on top of the
22-table 3NF base:
\begin{itemize}
  \item Composite index
    \texttt{idx\_audit\_event\_host\_timestamp} on
    \texttt{audit\_event(host\_id, timestamp)}, the index that powers
    the incident-response speedup documented in
    Section~\ref{sec:explain}.
  \item Saved view \texttt{v\_privilege\_escalation\_timeline}, an
    eight-table cross-domain join filtered to
    \texttt{attack\_phase.phase\_name = 'privilege\_escalation'} and
    aggregated with \texttt{string\_agg}, used by both the dashboard
    and the deep-dive query in Section~\ref{sec:view}.
\end{itemize}
A teammate can verify that both objects are present with one psql
command:

\begin{lstlisting}[language=bash, basicstyle=\ttfamily\small, frame=single]
docker exec security-logs-dev psql -U security_logs_user \
  -d security_logs \
  -c "SELECT version_num FROM alembic_version;" \
  -c "SELECT indexname FROM pg_indexes
      WHERE tablename = 'audit_event'
      ORDER BY indexname;" \
  -c "SELECT viewname FROM pg_views
      WHERE schemaname = 'public';"
\end{lstlisting}

Expected output: \texttt{version\_num = 742e860d116f},
\texttt{idx\_audit\_event\_host\_timestamp} appears in the index list,
and \texttt{v\_privilege\_escalation\_timeline} appears in the view
list.

\subsection{Final Row Counts}
Table~\ref{tab:rowcounts} reports the final loaded row counts (verified
2026-05-01 against the loaded database). Total across all 22 tables is
roughly 439{,}998 rows. The largest single table is
\texttt{labeled\_line\_rule} (184{,}651 rows: a per-rule annotation on
labeled lines), closely followed by \texttt{labeled\_line\_label}
(184{,}517 rows: the line-to-label many-to-many junction).

\begin{table}[htbp]
\caption{Final loaded row counts by domain (verified 2026-05-01).}
\label{tab:rowcounts}
\centering
\begin{tabular}{lll}
\toprule
\textbf{Domain} & \textbf{Table} & \textbf{Rows} \\
\midrule
Host inventory   & host                       & 22       \\
                 & host\_group                & 63       \\
                 & host\_fqdn                 & 20       \\
                 & host\_ipv4                 & 24       \\
                 & host\_ipv6                 & 24       \\
                 & host\_log\_config          & 66       \\
                 & os\_release                & 2        \\
\midrule
Audit events     & audit\_event               & 3{,}048  \\
                 & audit\_message             & 2{,}614  \\
                 & audit\_pam\_event          & 2{,}055  \\
                 & audit\_service\_event      & 555      \\
                 & audit\_login\_event        & 410      \\
                 & audit\_user\_login\_event  & 3        \\
                 & audit\_user\_cmd\_event    & 1        \\
                 & audit\_syscall\_event      & 8        \\
                 & audit\_avc\_event          & 8        \\
                 & audit\_proctitle\_event    & 8        \\
\midrule
Attack labels    & labeled\_line              & 61{,}862 \\
                 & labeled\_line\_label       & 184{,}517 \\
                 & labeled\_line\_rule        & 184{,}651 \\
                 & attack\_label              & 22       \\
                 & attack\_phase              & 7        \\
\midrule
\textbf{Total}   &                            & \textbf{$\approx$439{,}998} \\
\bottomrule
\end{tabular}
\end{table}

\subsection{Integrity Verification}
After load we run a small set of integrity checks. The two most
load-bearing checks are the orphan-FK count on \texttt{audit\_event}
and the cross-domain consistency check between \texttt{labeled\_line}
and \texttt{audit\_event} for rows where the labeled log is
\texttt{audit.log}. Both return zero on the current database.

\begin{lstlisting}[language=SQL,basicstyle=\ttfamily\scriptsize]
-- Every audit_event references a real host (FK consistency).
SELECT COUNT(*) AS orphans
FROM audit_event ae
LEFT JOIN host h ON h.host_id = ae.host_id
WHERE h.host_id IS NULL;
-- => 0 orphans

-- Every labeled_line line_number that targets audit.log lands
-- inside its host's audit_event range.
SELECT COUNT(*) AS unmatched
FROM labeled_line ll
WHERE ll.source_log = 'audit.log'
  AND NOT EXISTS (
    SELECT 1 FROM audit_event ae
    JOIN host h ON h.host_id = ae.host_id
    WHERE h.host_key   = ll.source_host
      AND ae.line_number = ll.line_number
  );
-- => 0 unmatched

-- Subtype coverage: every audit_event has exactly one subtype row.
WITH subtype_sum AS (
  SELECT
    (SELECT COUNT(*) FROM audit_pam_event)
  + (SELECT COUNT(*) FROM audit_service_event)
  + (SELECT COUNT(*) FROM audit_user_login_event)
  + (SELECT COUNT(*) FROM audit_user_cmd_event)
  + (SELECT COUNT(*) FROM audit_login_event)
  + (SELECT COUNT(*) FROM audit_syscall_event)
  + (SELECT COUNT(*) FROM audit_avc_event)
  + (SELECT COUNT(*) FROM audit_proctitle_event)
    AS subtype_total
)
SELECT (SELECT COUNT(*) FROM audit_event) AS total_events,
       subtype_total AS sum_of_subtypes
FROM   subtype_sum;
-- => total_events = sum_of_subtypes = 3,048
\end{lstlisting}

The third check (subtype-coverage equality) is the practical proof
that the supertype/subtype decomposition described in
Section~\ref{sec:schema} is total and disjoint, which is the property
that lets the report claim ``every audit event is exactly one
subtype.''
